%====================================================%
% SEZIONE 1: Notazione scientifica
%====================================================%
\begin{question}[ID=1.1, subtitle={Conversione in notazione scientifica — base}, tags={notazione scientifica}, difficulty=facile]
Trasforma le seguenti misure in notazione scientifica:
\begin{center}\begin{minipage}{0.92\linewidth}
\begin{multicols}{2}
\begin{enumerate}
\item $0{,}0045\ \text{m} = \underline{\hspace{2.5cm}}$
\item $32000\ \text{kg} = \underline{\hspace{2.5cm}}$
\item $0{,}00000082\ \text{s} = \underline{\hspace{2.5cm}}$
\item $5\,600\,000\ \text{J} = \underline{\hspace{2.5cm}}$
\end{enumerate}
\end{multicols}
\end{minipage}\end{center}
\begin{solution}
\[
4{,}5\times10^{-3}\ \text{m};\quad 3{,}2\times10^{4}\ \text{kg};\quad 8{,}2\times10^{-7}\ \text{s};\quad 5{,}6\times10^{6}\ \text{J}.
\]
\end{solution}
\end{question}

\begin{question}[ID=1.2, subtitle={Notazione scientifica di grandezze fisiche}, tags={notazione scientifica}, difficulty=facile]
Scrivi in notazione scientifica i seguenti valori fisici:
\begin{center}\begin{minipage}{0.92\linewidth}
\begin{multicols}{2}
\begin{enumerate}
\item $0{,}091\ \text{A} = \underline{\hspace{2.5cm}}$
\item $7{,}2\times 10^{5}\ \text{m} = \underline{\hspace{2.5cm}}$
\item $0{,}00031\ \text{N} = \underline{\hspace{2.5cm}}$
\item $48\,000\,000\ \text{Hz} = \underline{\hspace{2.5cm}}$
\end{enumerate}
\end{multicols}
\end{minipage}\end{center}
\begin{solution}
\[
9{,}1\times10^{-2}\ \text{A};\quad 7{,}2\times10^{5}\ \text{m} \text{ (già corretto)};\quad 3{,}1\times10^{-4}\ \text{N};\quad 4{,}8\times10^{7}\ \text{Hz}.
\]
\end{solution}
\end{question}

\begin{question}[ID=1.3, subtitle={Notazione scientifica con unità variegate}, tags={notazione scientifica}, difficulty=media]
Converti in notazione scientifica le seguenti grandezze:
\begin{center}\begin{minipage}{0.92\linewidth}
\begin{multicols}{2}
\begin{enumerate}
\item $0{,}00012\ \text{C} = \underline{\hspace{2.5cm}}$
\item $9\,300\,000\ \text{Pa} = \underline{\hspace{2.5cm}}$
\item $0{,}0084\ \text{m/s} = \underline{\hspace{2.5cm}}$
\item $6{,}75\times 10^{-4}\ \text{F} = \underline{\hspace{2.5cm}}$
\end{enumerate}
\end{multicols}
\end{minipage}\end{center}
\begin{solution}
\[
1{,}2\times10^{-4}\ \text{C};\quad 9{,}3\times10^{6}\ \text{Pa};\quad 8{,}4\times10^{-3}\ \text{m/s};\quad 6{,}75\times10^{-4}\ \text{F} \text{ (già corretto)}.
\]
\end{solution}
\end{question}


%====================================================%
% SEZIONE 2: Conversioni
%====================================================%
\begin{question}[ID=2.1, subtitle={Equivalenze di unità di misura — base}, tags={conversioni, unità di misura}, difficulty=facile]
Esegui le seguenti equivalenze:
\begin{center}\begin{minipage}{0.92\linewidth}
\begin{multicols}{2}
\begin{enumerate}
\item $0{,}045\ \text{km} = \underline{\hspace{2cm}}\ \text{m}$
\item $3200\ \text{g} = \underline{\hspace{2cm}}\ \text{kg}$
\item $0{,}008\ \text{A} = \underline{\hspace{2cm}}\ \text{µA}$
\item $5{,}2\ \text{kN} = \underline{\hspace{2cm}}\ \text{N}$
\item $3{,}5\ \text{mg} = \underline{\hspace{2cm}}\ \text{g}$
\item $2{,}4\ \text{km} = \underline{\hspace{2cm}}\ \text{mm}$
\item $0{,}0021\ \text{m}^3 = \underline{\hspace{2cm}}\ \text{cm}^3$
\item $7800\ \text{ns} = \underline{\hspace{2cm}}\ \text{s}$
\end{enumerate}
\end{multicols}
\end{minipage}\end{center}
\begin{solution}
45 m; 3,2 kg; 8000 µA; 5200 N; 0,0035 g; 2\,400\,000 mm; 2100 cm$^3$; $7{,}8\times10^{-6}$ s.
\end{solution}
\end{question}

\begin{question}[ID=2.2, subtitle={Conversioni di unità miste}, tags={conversioni, unità di misura}, difficulty=media]
Riscrivi nelle unità richieste:
\begin{center}\begin{minipage}{0.92\linewidth}
\begin{multicols}{2}
\begin{enumerate}
\item $3\ \text{s} = \underline{\hspace{2cm}}\ \text{ms}$
\item $6{,}4\ \text{MJ} = \underline{\hspace{2cm}}\ \text{J}$
\item $540000\ \text{mm} = \underline{\hspace{2cm}}\ \text{m}$
\item $1\ \text{m}^3 = \underline{\hspace{2cm}}\ \text{cm}^3$
\item $8{,}4\ \text{F} = \underline{\hspace{2cm}}\ \text{µF}$
\item $7{,}2\ \text{kmol} = \underline{\hspace{2cm}}\ \text{mol}$
\end{enumerate}
\end{multicols}
\end{minipage}\end{center}
\begin{solution}
3000 ms; $6{,}4\times10^{6}$ J; 540 m; $10^{6}$ cm$^3$; $8{,}4\times10^{6}$ µF; 7200 mol.
\end{solution}
\end{question}


%====================================================%
% SEZIONE 3: Sensibilità
%====================================================%
\begin{question}[ID=3.1, subtitle={Misure scorrette per sensibilità dello strumento}, tags={sensibilità, errore}, difficulty=facile]
Quale tra queste misure effettuate con uno strumento non può essere corretta:
\begin{center}\begin{minipage}{0.92\linewidth}
\begin{multicols}{2}
\begin{enumerate}[leftmargin=3.2cm]
\item $\ (1{,}05 \pm 0{,}05)\ mm$
\item $\ (2{,}17 \pm 0{,}05)\ mm$
\item $\ (10{,}2 \pm 0{,}1)\ mm$
\item $\ (10{,}21 \pm 0{,}1)\ mm$
\item $\ (11 \pm 1)\ mm$
\item $\ (11{,}1 \pm 1)\ mm$
\end{enumerate}
\end{multicols}
\end{minipage}\end{center}
\begin{solution}
Le misure scorrette sono \textbf{d)} e \textbf{f)}: l'incertezza deve avere lo stesso numero di cifre decimali della misura. In d) la misura ha 2 decimali ma l'incertezza ne ha 1; in f) la misura ha 1 decimale ma l'incertezza è intera.
\end{solution}
\end{question}

%====================================================%
% SEZIONE 4: Errore relativo, percentuale
%====================================================%
\begin{question}[ID=4.1, subtitle={Errore relativo e percentuale — confronto precisione}, tags={errori, relativo e percentuale}, difficulty=media]
Per ogni misura determina errore relativo e percentuale. Quale è più precisa?
\begin{center}\begin{minipage}{0.92\linewidth}
\begin{multicols}{2}
\begin{enumerate}[label=\alph*),leftmargin=1.8cm]
\item $(12{,}4 \pm 0{,}2)\ \text{m}$
\item $(8{,}20 \pm 0{,}05)\ \text{m}$
\item $(0{,}540 \pm 0{,}003)\ \text{m}$
\item $(1450 \pm 20)\ \text{m}$
\end{enumerate}
\end{multicols}
\end{minipage}\end{center}
\begin{solution}
\begin{align*}
\text{a)}\; &\varepsilon_r = 0{,}2/12{,}4 \approx 1{,}6\% &
\text{b)}\; &\varepsilon_r = 0{,}05/8{,}20 \approx 0{,}61\% \\
\text{c)}\; &\varepsilon_r = 0{,}003/0{,}540 \approx 0{,}56\% &
\text{d)}\; &\varepsilon_r = 20/1450 \approx 1{,}38\%
\end{align*}
La più precisa è \textbf{c)}, con l'errore relativo minore ($\approx 0{,}56\%$).
\end{solution}
\end{question}

\begin{question}[ID=4.2, subtitle={Errore relativo su corrente, lunghezza, mol, temperatura}, tags={errori, misure}, difficulty=media]
Calcola l'errore relativo, percentuale e indica la misura più precisa.
\begin{center}\begin{minipage}{0.92\linewidth}
\begin{multicols}{2}
\begin{enumerate}[label=\alph*),leftmargin=1.8cm]
\item $(3{,}50 \pm 0{,}04)\ \text{A}$
\item $(27{,}0 \pm 0{,}5)\ \text{cm}$
\item $(0{,}0021 \pm 0{,}0001)\ \text{mol}$
\item $(98{,}0 \pm 0{,}8)\ ^\circ\text{C}$
\end{enumerate}
\end{multicols}
\end{minipage}\end{center}
\begin{solution}
\begin{align*}
\text{a)}\; &\varepsilon_r \approx 1{,}1\% &
\text{b)}\; &\varepsilon_r \approx 1{,}9\% \\
\text{c)}\; &\varepsilon_r \approx 4{,}8\% &
\text{d)}\; &\varepsilon_r \approx 0{,}82\%
\end{align*}
La più precisa è \textbf{d)}.
\end{solution}
\end{question}

\begin{question}[ID=4.3, subtitle={Errore relativo su tensione, volume, massa, distanza}, tags={errori, misure}, difficulty=media]
Determina errore relativo, percentuale e quale misura risulta più precisa.
\begin{center}\begin{minipage}{0.92\linewidth}
\begin{multicols}{2}
\begin{enumerate}[label=\alph*),leftmargin=1.8cm]
\item $(4{,}80 \pm 0{,}02)\ \text{V}$
\item $(0{,}750 \pm 0{,}010)\ \text{L}$
\item $(120 \pm 3)\ \text{g}$
\item $(15{,}2 \pm 0{,}3)\ \text{m}$
\end{enumerate}
\end{multicols}
\end{minipage}\end{center}
\begin{solution}
\begin{align*}
\text{a)}\; &\varepsilon_r \approx 0{,}42\% &
\text{b)}\; &\varepsilon_r \approx 1{,}3\% \\
\text{c)}\; &\varepsilon_r \approx 2{,}5\% &
\text{d)}\; &\varepsilon_r \approx 2{,}0\%
\end{align*}
La più precisa è \textbf{a)}.
\end{solution}
\end{question}

\begin{question}[ID=4.4, subtitle={Errore relativo su massa, lunghezza, tempo, velocità}, tags={errori, misure}, difficulty=media]
Per ogni misura calcola errore relativo, percentuale e individua la più precisa.
\begin{center}\begin{minipage}{0.92\linewidth}
\begin{multicols}{2}
\begin{enumerate}[label=\alph*),leftmargin=1.8cm]
\item $(2{,}30 \pm 0{,}01)\ \text{kg}$
\item $(480 \pm 10)\ \text{mm}$
\item $(0{,}036 \pm 0{,}001)\ \text{s}$
\item $(9{,}0 \pm 0{,}4)\ \text{m/s}$
\end{enumerate}
\end{multicols}
\end{minipage}\end{center}
\begin{solution}
\begin{align*}
\text{a)}\; &\varepsilon_r \approx 0{,}43\% &
\text{b)}\; &\varepsilon_r \approx 2{,}1\% \\
\text{c)}\; &\varepsilon_r \approx 2{,}8\% &
\text{d)}\; &\varepsilon_r \approx 4{,}4\%
\end{align*}
La più precisa è \textbf{a)}.
\end{solution}
\end{question}

\begin{question}[ID=4.5, subtitle={Errore relativo su massa, distanza, corrente, temperatura}, tags={errori, misure}, difficulty=media]
Completa: errore relativo, errore percentuale e indica la misura più precisa.
\begin{center}\begin{minipage}{0.92\linewidth}
\begin{multicols}{2}
\begin{enumerate}[label=\alph*),leftmargin=1.8cm]
\item $(720 \pm 8)\ \text{kg}$
\item $(0{,}0150 \pm 0{,}0002)\ \text{m}$
\item $(6{,}4 \pm 0{,}2)\ \text{A}$
\item $(310 \pm 5)\ \text{K}$
\end{enumerate}
\end{multicols}
\end{minipage}\end{center}
\begin{solution}
\begin{align*}
\text{a)}\; &\varepsilon_r \approx 1{,}1\% &
\text{b)}\; &\varepsilon_r \approx 1{,}3\% \\
\text{c)}\; &\varepsilon_r \approx 3{,}1\% &
\text{d)}\; &\varepsilon_r \approx 1{,}6\%
\end{align*}
La più precisa è \textbf{a)}.
\end{solution}
\end{question}
